\item Cuando el aire se dilata adiabáticamente (sin ganar o perder calor), su presión $P$ y volumen $V$ están relacionados mediante la ecuación $PV^{1.4} = C$, en donde $C$ es una constante. Supóngase que en cierto instante el volumen es igual a $400 \text{ cm}^3$ y la presión es de $80 \text{ kPa}$ y ésta disminuye a razón de $10 \text{ kPa/min}$. ¿A qué velocidad está aumentando el volumen en dicho instante?
